% Related work (owner: B). Every claim about prior work is taken from the cited paper's abstract or
% from our notes in AGENTS.md; sentences on our own data say "our runs" or "our measurements".
% Entries marked TODO in references.bib still need their full metadata checked.

\section{Related Work}
\label{sec:related}

Our work sits between three lines of research: fingerprinting web agents from their traffic,
identifying a language model from its text, and tracking or linking users. Prior work studies each
on its own and mostly for browser agents. We combine them for coding agents, which fetch pages from
the user's own machine, and we target the user behind the agent rather than the tool.

\subsection{Fingerprinting AI agents from their traffic}

Borysenko~\cite{borysenko2026agents} measured nine coding agents visiting documentation portals and
found stable signatures in their User-Agent, headers and pre-fetch behaviour, with one to four
requests per prompt. Most agents use generic HTTP runtimes (Node's axios or got, curl, Go's net/http
or Colly), and only two name themselves. That study logged client IP addresses but did not analyse
them, and it did not include Codex CLI.

Kang et al.~\cite{kang2026whose} fingerprint six autonomous web agents across the TLS ClientHello
(JA3 and JA4), the HTTP/2 fingerprint, headers and browser events, and reach 97\% accuracy with an
Extra Trees classifier when the layers are combined; no single layer suffices. They explicitly
exclude coding agents, deliberately avoid logging participant IP addresses, and their goal is content
protection, so linking requests to a user is left open. Our logging server follows their layered
approach, and our harness classifier uses the same model family.

Wang et al.~\cite{wang2026fpagent} compare seven browsing agents with humans and show that local
agents expose the host machine's attributes while cloud agents expose a fixed remote environment.
They also observe that agents produce mouse events only at click time and paste rather than type, so
human behavioural biometrics do not transfer. Lugoloobi et al.~\cite{lugoloobi2026actions} identify
which of 14 models drives a browser agent from its UI traces, but this needs rich interaction that
coding agents do not produce. Fayolle et al.~\cite{fayolle2026browser} send six web agents to
honeysites protected by anti-bot mechanisms (robots.txt, CAPTCHAs, proof-of-work and Cloudflare's
free tier). Some agents bypass all of them, yet network-, HTTP- and browser-level fingerprints tell
every agent apart from humans and from each other, and stealth features often make agents easier to
detect. Unlike all of these, we study coding and command-line agents, which make only a handful of
plain HTTP requests and run no JavaScript.

\paragraph{Active detection.} A site can also probe the visitor. The LLM agent honeypot of
Palisade~\cite{reworr2024honeypot} combines prompt injection with timing to detect LLM-driven
attackers, and \emph{Say Shibboleth}~\cite{shibboleth2026} uses traps. We instead measure what a site
learns \emph{passively}, from traffic it receives anyway; our input page asks for a summary only in
plain visible text, never through hidden instructions.

\subsection{Identifying the model from its text}

LLMmap~\cite{pasquini2025llmmap} identifies the model behind an application from its responses,
across 42 model versions. Sun et al.~\cite{sun2025idiosyncrasies} show that a classifier on top of
an LLM2Vec text encoder~\cite{behnamghader2024llm2vec} tells apart the outputs of five chatbots with
high accuracy. Our model track reuses their published classifier as a frozen backbone and trains new
output heads on our current models, applied to the short text an agent submits on our page. We
further show that, for one agent, the sequence of tools it uses already reveals the model tier
without any text.

\subsection{Tracking and linking users}

Cross-device tracking links a user's devices from IP addresses and browsing history
(Zimmeck et al.~\cite{zimmeck2017crossdevice}, F1 0.91), and browsing histories alone re-identify a
large share of users (Bird et al.~\cite{bird2020replication}). Crichton et al.~\cite{crichton2025rethinking}
show, over 150,000 users and two years, that behavioural fingerprints are highly unique and stable,
and that combining them with browser fingerprints works best. These are our baselines for linking.

UID smuggling (Randall et al.~\cite{randall2022uid}) passes an identifier between sites through
navigation URLs. We observe an agent variant of it: the agent relays the visited URL, including our
per-visit identifier, to the user, whose own browser may then carry it. In 5 of our 27 Claude Code
runs the agent's answer included the visited URL, all of them WebFetch runs.

Two older linking signals are TCP timestamp clock skew~\cite{kohno2005clock} and TLS session
tickets~\cite{sy2018tls}. In our measurements neither helps in our setting: the TCP timestamps were
randomised per connection on current Linux and absent by default on Windows, and none of our agent
connections resumed a TLS session, so no ticket links them.
% TODO: point to our results section for these two measurements once it exists.

Zhang et al.~\cite{zhang2025traffic} infer agent behaviour and some user attributes from encrypted
traffic, but as an on-path network observer rather than the visited website. Ukani
et al.~\cite{ukani2025privacy} show that several browser agents keep state across sessions and leak
details from earlier chats to websites, without studying fingerprinting or linking.

\paragraph{Our position.} To our knowledge, no prior work links the requests of tool-using coding
agents to the person who runs them, across agents, partner sites and the person's own browser, from
traffic a site receives passively.
